\section{Design Summary}
To map the bigger picture of the design of this application, here are the top level diagrams for the application UI flow, UI state and concurrency diagrams

\section{Application Architecture}
\label{sec:app-architecture}

The application comprises four concurrent threads — RAG, Summarisation, UI, and Transcription — each responsible for a distinct subsystem. At program start the UI thread spawns the three worker threads, each of which loads the models required for its subsystem. Figure~\ref{fig:concurrency-diagram} shows the interactions between threads across the four phases of the application lifecycle.

\begin{landscape}
\begin{figure}[H]
\centering
\resizebox{\linewidth}{!}{%
\begin{tikzpicture}[
    >=Stealth,
    event/.style={circle, fill=black, inner sep=0pt, minimum size=6pt},
    arr/.style={->, semithick},
    x=1.0cm, y=1.0cm,
]

%% Lane y-coordinates (centre of each 2.5-unit band)
\def\yRAG{8.75}
\def\ySumm{6.25}
\def\yUI{3.75}
\def\yTrans{1.25}

%% Lane boundary lines  (total width ~42)
\foreach \y in {0, 2.5, 5, 7.5, 10}
    \draw[cyan!70!blue, thick] (0, \y) -- (42.0, \y);

%% Lane labels
\node[anchor=east, font=\Large, align=right] at (-0.2, \yRAG)   {RAG\\thread};
\node[anchor=east, font=\Large, align=right] at (-0.2, \ySumm)  {Summarisation\\thread};
\node[anchor=east, font=\Large, align=right] at (-0.2, \yUI)    {UI\\thread};
\node[anchor=east, font=\Large, align=right] at (-0.2, \yTrans) {Transcription\\thread};

%% Phase separator lines
\foreach \x in {11.0, 21.0, 33.0}
    \draw[densely dashed, gray!60] (\x, 10.7) -- (\x, -0.5);

%% Phase labels
\node[font=\large\itshape] at (5.5,  11.1) {Program start};
\node[font=\large\itshape] at (16.0, 11.1) {Start recording};
\node[font=\large\itshape] at (27.0, 11.1) {Stop recording};
\node[font=\large\itshape] at (37.5, 11.1) {View guidance};

%%====================================================================
%% PHASE 1 — Program start: UI spawns the three worker threads
%%====================================================================

\node[event] (U0)  at (1.5, \yUI)    {};

%% UI → RAG
\node[event] (R0)  at (3.5, \yRAG)   {};
\node[font=\small, above=4pt] at (R0) {load MedCPT (NPU)};
\draw[arr] (U0) -- (R0);

%% UI → Summarisation
\node[event] (S0)  at (3.5, \ySumm)  {};
\node[font=\small, above right=4pt, align=left] at (S0) {load Med42-8B (GPU)};
\draw[arr] (U0) -- (S0);

%% UI → Transcription
\node[event] (Tr0) at (3.5, \yTrans) {};
\node[font=\small, below=4pt, align=center] at (Tr0)
    {load ECAPA-TDNN\\+ Whisper (CPU)};
\draw[arr] (U0) -- (Tr0);

%% UI → UI View
\node[event] (V_loading) at (3.5, \yUI) {};
\node[font=\small, above=4pt, align=center] at (V_loading)
    {View:\\loading application};
\draw[arr] (U0) -- (V_loading);
%% Transcription: models finished loading
\node[event] (Tr_ready) at (7.5, \yTrans) {};
\node[font=\small, below=4pt] at (Tr_ready) {models loaded};
\draw (Tr0) -- (Tr_ready);

%% Transcription → UI: signals models loaded
\node[event] (U_modload) at (8.0, \yUI) {};
\node[font=\small, above=4pt, align=center] at (U_modload) {View:\\ready to record};
\draw[arr] (Tr_ready) -- (U_modload);
\draw[arr] (V_loading) -- (U_modload);


%%====================================================================
%% PHASE 2 — Start recording
%%====================================================================

%% UI: user presses start recording → signals Transcription directly
\node[event] (U_srec) at (14.0, \yUI) {};
\node[font=\small, above=4pt, align=center] at (U_srec) {start recording\\pressed};
\draw[arr] (U_modload) -- (U_srec);

\node[event] (V_recording) at (17.5, \yUI) {};
\draw[arr] (U_srec) -- (V_recording);
\node[font=\small, above=4pt, align=center] at (V_recording) {View:\\recording};

%% Transcription: receives signal from UI, begins transcribing
\node[event] (Tr_sig) at (17.5, \yTrans) {};
\draw[arr] (U_srec) -- (Tr_sig);
\node[font=\small, below=4pt] at (Tr_sig) {start transcription};
\draw[arr] (Tr_ready) -- (Tr_sig); %% horizontal: models loaded → start transcription

%%====================================================================
%% PHASE 3 — Stop recording
%%====================================================================

%% UI: stop recording pressed → signals Transcription
\node[event] (U_stop) at (21.5, \yUI) {};
\node[font=\small, above=4pt, align=center] at (U_stop) {stop recording\\pressed};
\draw[arr] (V_recording) -- (U_stop); %% horizontal: View: recording → stop recording pressed

%% Transcription: receives stop signal, sends back accumulated buffer
\node[event] (Tr_send) at (25.0, \yTrans) {};
\draw[arr] (U_stop) -- (Tr_send);
\node[font=\small, below=4pt, align=center] at (Tr_send)
    {complete \& send\\transcription};
\draw[arr] (Tr_sig) -- (Tr_send);

%% UI: show the generating summarisation view
\node[event] (V_gensum) at (25.0, \yUI) {};
\draw[arr] (U_stop) -- (V_gensum);
\node[font=\small, above=4pt, align=center] at (V_gensum)
    {View:\\generating\\summarisation};

%% Transcription → Summ: dispatch transcription directly
\node[event] (S_stop) at (28.5, \ySumm) {};
\node[font=\small, above=4pt] at (S_stop) {summarise};
\draw[arr] (Tr_send) -- (S_stop)
    node[pos=0.5, below right=3pt, font=\small] {dispatch transcription};
\draw[arr] (S0) -- (S_stop); %% horizontal: load Med42-8B → summarise

%% Summ: summarisation complete
\node[event] (S_summ) at (31.0, \ySumm) {};
\node[font=\small, above=4pt] at (S_summ) {summary complete};
\draw (S_stop) -- (S_summ);

%% Summ → UI: sends summary back to UI
\node[event] (U_recv) at (32.5, \yUI) {};
\node[font=\small, below=4pt, align=center] at (U_recv) {View:\\summarisation};
\draw[arr] (S_summ) -- (U_recv);

%% V_gensum → U_recv: UI transitions directly to summarisation view
\draw[arr] (V_gensum) -- (U_recv);

%%====================================================================
%% PHASE 4 — View guidance
%%====================================================================

%% UI: user presses view guidance → goes directly to RAG
\node[event] (U_view) at (35.5, \yUI) {};
\node[font=\small, below=4pt, align=center] at (U_view)
    {view guidance\\pressed};
\draw[arr] (U_recv) -- (U_view);
\draw[arr] (S_summ) -- (U_recv);

%% RAG: receives request directly from UI
\node[event] (R_embed) at (36.5, \yRAG) {};
\node[font=\small, above=4pt] at (R_embed) {embed/lookup};
\draw[arr] (U_view) -- (R_embed);
\draw[arr] (R0) -- (R_embed); %% horizontal: load MedCPT → embed/lookup

%% RAG: performs lookup, result ready
\node[event] (R_done) at (39.5, \yRAG) {};
\node[font=\small, above=4pt] at (R_done) {send guidance};
\draw[arr] (R_embed) -- (R_done);

%% RAG → UI: returns found guidance
\node[event] (U_done) at (41.0, \yUI) {};
\node[font=\small, below=4pt, align=center] at (U_done) {View:\\medical };
\draw[arr] (R_done) -- (U_done);

\end{tikzpicture}%
}
\caption{Concurrency diagram illustrating thread interactions across the four phases of the application lifecycle.}
\label{fig:concurrency-diagram}
\end{figure}
\end{landscape}
